\documentclass[12pt,a4paper]{article}

\usepackage[utf8]{inputenc}
\usepackage[T1]{fontenc}
\usepackage[brazil]{babel}
\usepackage[top=3cm,left=3cm,bottom=2cm,right=2cm]{geometry}
\usepackage{setspace}
\usepackage{graphicx}
\usepackage[alf]{abntex2cite}
\onehalfspacing

\title{$titulo}
\author{$autor \\ Nome da instituição}
\date{}

\begin{document}

\maketitle

\begin{abstract}
Entre 100 e 250 palavras, em um parágrafo só: objetivo, método, resultados e conclusões. Logo abaixo, de três a cinco palavras-chave.

\textbf{Palavras-chave}: primeira; segunda; terceira.
\end{abstract}

\section{Introdução}

Contextualize o tema, apresente o problema de pesquisa, a justificativa e os objetivos. Na ABNT, as citações aparecem com o sobrenome do autor e o ano, como em \cite{barbosa2010ihc}, ou no meio da frase, como fazem \citeonline{barbosa2010ihc}.

\section{Referencial Teórico}

Apresente os conceitos e autores em que o trabalho se apoia.

\section{Metodologia}

Classifique a pesquisa (por exemplo, qualitativa e exploratória) e descreva os procedimentos e instrumentos usados.

\section{Resultados e Discussão}

Apresente e interprete os resultados, comparando com o referencial teórico.

\section{Considerações Finais}

Responda ao objetivo, aponte as limitações e sugira trabalhos futuros.

\bibliographystyle{abntex2-alf}
\bibliography{referencias}

\end{document}
